Indoor percussion is a judged music-and-movement activity in which ensembles
perform designed shows combining percussion music with staged visual movement.
Adjudicated scores are expected to rise as groups refine these shows across a
season.
We study whether an ensemble's debut, its first valid non-championship
competitive appearance, can predict its terminal championship subtotal, the
combined caption score (sum of adjudicated scoring dimensions) before timing penalties.
Using a newly engineered database of 4,441 Southern California Percussion
Alliance (SCPA) performances spanning nine seasons (2017--2026), we construct
a cohort of 781 ensemble-seasons, where each observation represents one
performing ensemble in one season. Each has both a debut and a championship
result. Models are fit on 574 development observations and
evaluated on 207 held-out observations from 2025--2026. Against three baselines (global
median, class-level median, and debut-score identity), we evaluate Ridge
regression and gradient boosting. Ridge achieves a held-out mean absolute
error of 2.676 points (RMSE 3.499, $R^2 = 0.769$, Spearman $\rho = 0.892$),
substantially outperforming all baselines. This result reflects predictability
within the recorded adjudication system: the model forecasts a later
adjudicated score from an earlier one and does not measure or explain the
underlying performance quality, artistry, or season development that produced
either score. Findings are further limited by championship-selection bias,
observational judge assignments, and SCPA-only coverage.
